\subsection{Семинары «SOLID» (8 очных часов)}

В течение четырёх пар практических занятий предполагается внедрить принципы SOLID
в код, полученный в первой практической работе, одновременно улучшив отзывчивость
консольного интерфейса (CLI).

В этой работе (и далее в курсе) мы будем придерживаться очень строгого
(и даже оверинжинирингового) подхода к реализации SOLID: будем выполнять
принцип S максимально строго, даже если в реальной практике так делать не стоит.

Порядок вашей работы:

\begin{itemize}
    \item реализовать отзывчивый CLI (то есть сделать максимально удобными ввод и вывод
    значений во всех ситуациях);
    \item реализовать принцип S так, чтобы каждый класс имел строго одну ответственность
    в максимально строгом смысле (по мере демонстрации работ преподавателю вы будете
    получать указания, что ещё следует выделить);
    \item реализовать принцип D так, чтобы почти каждый класс реализовывал интерфейс
    (мы будем создавать абстрактные классы без свойств и только с абстрактными методами
    и использовать множественное наследование); если необходимо использовать объекты
    других классов, передавайте их в конструктор или метод не как конкретные типы,
    а как типы интерфейсов;
    \item реализовать принцип I так, чтобы в конструкторы и методы передавались
    максимально узкие интерфейсы;
    \item реализовать принцип O так, чтобы класс был открыт для расширения,
    но закрыт для изменения;
    \item считать, что принцип L реализуется автоматически, если вы следуете
    правилам наследования, изложенным на первом занятии;
    \item реализовать автоматизированное модульное (unit) тестирование всех
    полученных классов (для этого понадобится создать mock-объекты или mock-классы);
    \item добавить документацию (docstring).
\end{itemize}